\BourbakiUnnumberedChapter{Historical Note}{Historical Note (Chapters IV and V)}

(Chapters IV and V)

(Numbers in brackets refer to the bibliography at the end of this Note.)

The theory of commutative fields --- and the theory of polynomials which is closely related to it --- derives directly from what constituted until the middle of the 19th century the principal object of classical Algebra: the solution of algebraic equations, and of the problems of geometrical constructions equivalent to them.

Once one seeks to solve an algebraic equation of degree \(>1\) , one finds oneself in the presence of quite new computational difficulties, as it is no longer possible to determine the unknown by performing « rational » calculations on the data. This difficulty must have been perceived at a very early stage; it must be reckoned as one of the most important contributions of the Babylonians that they were able to reduce the solution of quadratic and biquadratic equations to a single new algebraic operation, the extraction of square roots (as is shown by the many numerical equations which have been found solved in this way in the texts that have come down to us ({[}1{]}, p.~183-193)). As far as the formal calculation is concerned, Antiquity never passes beyond this point in the problem of the solution of algebraic equations; the Greeks of the classical epoch in effect confine themselves to restating the Babylonian formulae in geometrical terms, and their use in algebraic form is not documented before Hero (100 AD) and Diophantus.

It is in quite a different direction that the Greeks made decisive progress. We have little information about the way in which the Babylonians conceived and calculated the square roots of integers which were not perfect squares *; in the rare texts on this question which have come down to us, they seem to have contented themselves with fairly coarse methods of approximation ({[}1{]}, p.~33-38). The Pythagorean School, which had quite rigorously determined the concept of commensurable magnitudes and attached to it a quasi-religious character, was not able to keep to this point of view; possibly it was the failure of repeated attempts to express \(\sqrt{2}\) rationally that led them finally to demonstrate that this number is irrational **.

We have said elsewhere (cf.~Hist. Note to Book III, Chapter IV) how this discovery, which marks a watershed in the history of Mathematics, reacted profoundly on the conception of « number » among the Greeks, and led them to create an algebra of exclusively geometrical character, in order to find a mode of representation (or perhaps an « existence » proof) for the incommensurable ratios, which they refused to consider as numbers. Most frequently an algebraic problem was reduced by them to the intersection of two auxiliary plane curves, appropriately chosen, or to several successive determinations of such intersections. Late traditions of little authority ascribe to Plato the introduction of a first classification of these constructions, destined to a long and brilliant career : for reasons that are philosophical rather than mathematical, it seems, that what are known as « ruler-and-compass-constructions », that is those where the only auxiliary curves occurring are straight lines and circles, have been put on one side *. In any case Euclid in his Elements {[}2{]} limits himself exclusively to treating problems that can be solved in this fashion (without however characterizing them by a particular name) ; a circumstance which undoubtedly contributed no little to fix the attention of mathematicians of succeeding centuries on these problems. But we know today ** that the algebraic equations that may be solved « by ruler and compasses » are of a very special type ; in particular, an irreducible equation

leads immediately to a demonstration of the irrationality of \(\sqrt{5}\) , and he has put forward the hypothesis (which unfortunately is not supported by any text) that it was in this fashion that the Pythagoreans discovered irrational numbers (K. von FRITZ, Ann. of. Math., vol.~46 (1945), p.~242).

* In connection with this principle one also attributes to Plato the classification of plane curves into « plane loci » (straight line and circle), « solid loci » (the conics, obtained by plane section of a solid body, the cone), all the other curves being grouped together under the name « τόποι γραμμικοι ». It is curious to see the influence of this classification extending even to Descartes, who in his Geometry arranges in the same « genus » the equations of degree 2n -- 1 and 2n, no doubt because those of degree 1 or 2 are solved by intersections of « plane loci » and those of degree 3 or 4 by intersections of « solid loci ».

** The determination of the points of intersection of a straight line and a circle (or of two circles) is equivalent to finding the solution of an equation of the second degree whose coefficients are rational functions of the coefficients of the equations of the straight line and circle (or the two circles) considered. It follows easily that the coordinates of a point constructed « by ruler and compasses » from given points belong to an extension L of the field Q of rational numbers, obtained as follows : if K is the field obtained by adjoining to Q the coordinates of the given points, then there exists an increasing sequence \((\mathrm{L}_{i})_{0\leq i\leq n}\) of fields intermediate between K and L, satisfying the conditions \(K=L_{0}\) , \(L=L_{n}\) , \([L_{i}:L_{i-1}]=2\) for \(1\leq i\leq n\) . By induction on n we deduce that the degree over K of the Galois extension N generated by L is a power of 2; conversely it may be shown that if this condition holds, then there exists a sequence \((\mathrm{L}_{i})\) of fields intermediate between K and L with the above properties, and therefore the problem posed can be solved by ruler and compasses (cf.~N. TSCHEBOTARÖW, Grundzüge der Galoisschen Theorie (transl. H. Schwerdtfeger), Groningen (P. Noordhoff), 1950, p.~351).

of the 3rd degree (over the field of rational numbers) cannot be solved in this way, and the Greeks had encountered quite early problems of the 3rd degree which have remained famous, such as the duplication of the cube (solution of \(x^3 = 2\) ) and the trisection of an angle; on the other hand the quadrature of the circle placed them in the presence of a transcendental problem. We find that to solve these problems they introduce numerous curves, both algebraic (conics, cissoid of Diocles, conchoid of Nicomedes) and transcendental (quadratrix of Dinostratus, Archimedes' spiral); all this inevitably led them to conduct an independent study of these various curves, thus paving the way to the future developments of Analytic Geometry, Algebraic Geometry and of Infinitesimal Calculus. But these methods contribute no progress to the solution of algebraic equations *, and the only work of Antiquity which has made a notable contribution to this question and exerted a lasting influence on the algebraists of the Middle Ages and the Renaissance is Book X of Euclid's Elements {[}2{]}; in this book (of which the principal results are ascribed by certain historians to Theaetetus), he considers the expressions obtained by the combination of several radicals such as \(\sqrt{\sqrt{a} \pm \sqrt{b}}\) ( \(a\) and \(b\) rational), gives conditions under which these expressions are irrational, classifies them in numerous categories (which he shows to be distinct) and studies the algebraic relations between these various irrationals, such as those which we would write today as

\[
\sqrt {\sqrt {p} + \sqrt {q}} = \sqrt {\frac {1}{2} (\sqrt {p} + \sqrt {p - q})} + \sqrt {\frac {1}{2} (\sqrt {p} - \sqrt {p - q})};
\]

all expressed in the usual geometric language of the Elements, which makes the exposition particularly dense and awkward.

After the decline of classical Greek Mathematics the notions relating to algebraic equations underwent modification. There can be no doubt that during the whole of the classical period the Greeks possessed methods of indefinite approximation of square roots, about which we are unfortunately ill informed **.

With the Hindus, later the Arabs and their Western emulators of the Middle Ages, the extraction of roots of all orders becomes an operation which tends to be considered as fundamental with the same right as the rational operations of algebra, and to be denoted, like the latter, by symbols which are more and more easily handled in calculations *. The theory of the equation of the second degree which was perfected during the Middle Ages (number of roots, negative roots, case of impossibility, double root), and that of biquadratic equations, provided models of formulae for the solution of equations « by radicals », on which algebraists will attempt for centuries to model analogous formulae for the solution of equations of higher degree, in the first instance for the equation of the 3rd degree. Leonardo of Pisa, who was chiefly responsible for introducing Arab science to the West in the 13th century, recognized that in any case the irrationals classified by Euclid in his tenth book cannot serve for this purpose (a new impossibility proof, in a theory which abounds with them), and we see him trying already the analogous calculations for cube roots, obtaining relations such as

\[
\sqrt [ 3 ]{1 6} + \sqrt [ 3 ]{5 4} = \sqrt [ 3 ]{2 5 0},
\]

analogous to Euclid's formulae for square roots (and of which we actually find earlier examples among the Arabs). But three centuries of fruitless efforts are still to pass before Scipio del Ferro, at the beginning of the 16th century, finally arrives at the formula for the solution of the equation \(x^3 + ax = b\) :

\[
x = \sqrt [ 3 ]{\frac {b}{2} + \sqrt {\left(\frac {b}{2}\right) ^ {2} + \left(\frac {a}{3}\right) ^ {3}}} + \sqrt [ 3 ]{\frac {b}{2} - \sqrt {\left(\frac {b}{2}\right) ^ {2} + \left(\frac {a}{3}\right) ^ {3}}}.\tag{1}
\]

We refrain from describing here the picturesque side of this sensational discovery --- the quarrels it provoked between Tartaglia on the one hand and Cardano and his school on the other --- as well as the often engaging figures of the scholars who were the protagonists. But we must note the decisive progress ensuing in the theory of equations at the hands of Cardano and his pupils. Cardano, who has less aversion than the majority of his contemporaries to the use of negative numbers, observes thus that the equations of third degree may have three roots, and the biquadratic equations four ({[}3{]} vol.~IV, p.~259), and he remarks that the sum of the three roots of \(x^3 + bx = ax^2 + c\) (an equation in which he knows how to make the term in \(x^2\) vanish) is always equal to \(a\) (ibid.). Guided no doubt by this relation and the intuition of its general character, he has the first idea of the notion of the multiplicity of a root; above all he is emboldened (not without oratory precautions), to calculate formally with expressions containing square roots of negative numbers. It is probable that he was led to this by the fact that such expressions arise naturally in the use of the formula (1) when \(\left(\frac{b}{2}\right)^{2} + \left(\frac{a}{3}\right)^{3} < 0\) (the so-called «irreducible case», in which Cardano recognized the existence of three real roots); in any case this is brought out clearly by his disciple R. Bombelli who in his Algebra ({[}4{]}, p.~293) proves the relation

\[
\sqrt [ 3 ]{2 + \sqrt {- 1 2 1}} = 2 + \sqrt {- 1}
\]

and who takes care to give explicitly the rules of calculating with complex numbers in a form which is already very close to modern expositions *. Finally in 1545 another student of Cardano, L. Ferrari, succeeds in solving the general equation of the 4th degree, by means of an auxiliary equation of the 3rd degree **.

After such rapid progress, the period that followed, until the middle of the 18th century, barely develops the new ideas introduced by the Italian School. Thanks to the essential progress which he brings to algebraic notation, Vieta is able to express in general fashion the relations between the coefficients and the roots of an algebraic equation, at least when the roots are all positive *** ({[}5{]}, p.~158). More daring, A. Girard {[}6{]} does not hesitate to affirm (of course without proof) that an equation of degree \(n\) has exactly \(n\) roots, provided that the « impossible roots » are counted, each with its degree of multiplicity, and that these roots satisfy the relations given by Vieta; he also obtains, for the first time, the expression for the sums of equal powers of the roots, up to exponent 4.

But the spirit of the 17th century is turned towards other directions and it is only indirectly that Algebra reaps some benefit from the new discoveries of Analytic Geometry and Infinitesimal Calculus. Thus Descartes' method of obtaining the

* Bombelli ({[}4{]}, p.~169 and 190) considers complex numbers as « linear combinations » with positive coefficients, of four basis elements: « piu » (+1), « meno » (-1), « piu de meno » (+i) and « meno de meno » (-i); in particular he postulates an axiom that « piu » and « piu de meno » cannot be added, the first appearance of the notion of linear independence.

** The equation being reduced to the form \(x^4 = ax^2 + bx + c\) , we determine a number \(z\) so that the right-hand side of the equation

\[
(x ^ {2} + z) ^ {2} = (a + 2 z) x ^ {2} + b x + (c + z ^ {2})
\]

is a perfect square, which gives an equation of the third degree for \(z\) .

*** Vieta, a passionate admirer of the Ancients, refrains systematically from introducing negative numbers into his arguments; he is nonetheless able on occasion to express in his language the relations between coefficients and roots when certain of the latter are negative; for example, if the equation \(x^3 + b = ax\) has two positive roots \(x_1, x_2\) ( \(a > 0, b > 0\) ), Vieta shows that \(x_1^2 + x_2^2 + x_1x_2 = a\) and \(x_1x_2(x_1 + x_2) = b\) ({[}5{]}, p.~106).

tangents to algebraic curves (cf.~Hist. Note to Book IV, Chap. I-II-III, p.~46) is related to the criterion for the multiplicity of a root of an algebraic equation, stated by his disciple Hudde ({[}7{]}, p.~433 and 507-509). It is no doubt also to Descartes' influence that the distinction between algebraic functions and transcendental functions should be ascribed, parallel to that which he introduces in his Géométrie between the « geometric » curves and the « mechanical » curves (cf.~Hist. Note to Book IV, Chap. I-II-III, p.~46 and 61). In any case this distinction is made perfectly clear by J. Gregory who in 1667, seeks even to prove that the area of a circular sector cannot be an algebraic function of the chord and the radius *. The expression « transcendental » is due to Leibniz, whom these questions of classification never cease to interest all through his career and who, about 1682 discovers a simple proof of the result Gregory was pursuing, by proving that sin x is not an algebraic function of x ({[}8{]}, vol.~V, p.~97-98) **. With his friend Tschirnhaus, Leibniz is actually one of the few mathematicians of his time to interest himself still in the problem of the solution « by radicals » of algebraic equations. At his beginnings we find him studying the « irreducible case » of the equation of the 3rd degree, and convince himself (as it happens, with insufficient proof) that it is impossible in this case to eliminate imaginary quantities from the formulae of the solution ({[}9{]}, p.~547-564). At about the same time he attacks, also without success, the solution by radicals of the equation of the 5th degree ; and when later Tschirnhaus claims to solve the problem by making all the terms of the equation disappear except the two extremes, by a transformation of the form y = P(x), where P is a suitably chosen polynomial of the 4th degree, Leibniz notices at once that the equations which determine the coefficients of P(x) are of degree \textgreater{} 5, and regards the method as doomed to failure ({[}9{]}, p.~402-403).

It seems to have been the needs of the new Analysis which gradually rekindled the interest in algebra. The integration of rational fractions, effected by Leibniz and Johann Bernoulli, and the question of imaginary logarithms closely related to it, provided an opportunity to develop the calculations with imaginary numbers in greater depth and to take up again the question of the decomposition of a polynomial into factors of the first degree (« fundamental theorem of algebra ») *. Right at the beginning of the 18th century, the solution of the binomial equation \(x^n - 1 = 0\) is reduced by Cotes and de Moivre to the division of the circle into \(n\) equal parts; to obtain the expressions of its roots « by radicals » it is therefore enough to know how to do it when \(n\) is an odd prime, and de Moivre remarks that the substitution \(y = x + \frac{1}{x}\) then reduces the problem to the solution « by radicals » of an equation of degree \((n - 1)/2\) . As far as the « fundamental theorem » is concerned, after the repeated failures of a general solution « by radicals » (including several attempts by Euler {[}10{]}), efforts are beginning to tend towards a priori proofs, not using explicit formulae of solution. Without entering into the details of the methods proposed (which were to bear fruit in the proofs of Lagrange and Gauss; cf.~Hist. Notes to Book II, Chap. VI-VII and Book III, Chap. VIII), it is appropriate to note here the point of view from which the problem was seen in the middle of the 18th century; it is admitted (without any justification other than a vague feeling of generality, coming no doubt, as with A. Girard, from the existence of relations between the coefficients and the roots) that an equation of degree \(n\) always has \(n\) « ideal » roots, an which calculations may be performed as on numbers, without knowing whether they are numbers (real or complex); what has to be proved (using if necessary calculations on the ideal roots), is that at least one of these roots is an ordinary complex number **. In this defective form the first germ of the general idea of « formal adjunction » may be recognized, which in spite of Gauss's objections ({[}13{]}, vol.~III, p.~1) was to become the basis of the modern theory of commutative fields.

With the fundamental memoirs of Lagrange ({[}11{]}, \(a\) ) and Vandermonde {[}12{]} the year 1770 sees the opening of a new and decisive period in the history of algebraic equations. The empiricism of more or less lucky attempts to find formulae of solution, which until then had reigned undivided, was followed by a systematic analysis of the problems posed and of the methods capable of solving them, an analysis which in sixty years led to the definitive results of Galois. Both Lagrange and Vandermonde started out from the ambiguity introduced by the multiple determinations of the radicals in the formulae for the solution of equations of degree \(\leqslant 4\) ; this fact had already attracted the attention of Euler ({[}10{]}, \(a\) ), who had shown among other things how in del Ferro's formula one has to associate the radicals occurring there in such a way as to obtain 3 roots and not 9. Lagrange remarked that each of the cube roots in del Ferro's formula may be written in the form \(\frac{1}{3} (x_1 + \omega x_2 + \omega^2 x_3)\) , where \(\omega\) is a cube root of unity and \(x_1, x_2, x_3\) the roots of the proposed equation, taken in a certain order, and he made the fundamental observation that the function \((x_1 + \omega x_2 + \omega^2 x_3)^3\) of the three roots can only take two distinct values for every permutation of the three roots, which explains a priori the success of the methods of solution of this equation. A similar analysis of the methods of solution of the equation of the 4th degree leads him to the function \(x_1 x_2 + x_3 x_4\) of the four roots, which only takes three distinct values for every permutation of the roots and is therefore root of an equation of the third degree with coefficients which are rational functions in those of the given equation *; these facts constitute as Lagrange says « the true principles and, so to speak, the metaphysics ** of the solution of equations of the 3th and 4th degree » ({[}11{]}, a), p.~357). Guided by these examples he plans to study in general, for an equation of degree \(n\) the number \(\nu\) of values *** which a rational function V of the \(n\) roots may assume when the latter are permuted arbitrarily; he thus in reality inaugurates (in this terminology still closely adapted to the theory of equations) the theory of groups and that of fields, of which he already obtains several fundamental results by making use of the same principles as those that are employed today. For example, he shows that the number \(\nu\) is a divisor of \(n!\) , by the reasoning which today serves to prove that the order of a subgroup of a finite group divides the order of this group. Still more remarkable is the theorem in which he shows that if \(V_1\) and \(V_2\) are two rational functions of the roots such that \(V_1\) and \(V_2\) remain invariant under the same permutations, then either is a rational function of the other and the coefficients of the equation (a particular case of the theorem of Galois characterizing a subextension of a Galois extension as field of invariants of its Galois group): « This problem » he says « seems to me one of the most important of the theory of equations, and the general solution we shall give will serve to throw new light on this part of Algebra » ({[}11{]}, a), p.~374).

All these researches are naturally, in Lagrange's mind, preliminaries to the analysis of possible methods of solution of algebraic equations by successive reduction to equations of lower degree, such a method being related, as he shows, to the formation of rational functions of the roots taking fewer than \(n\) values under permutation of the roots. Guided no doubt by results on the equation of the 3rd degree, he introduces in general the « Lagrange resolvents » \(y_{k} = \sum_{h=1}^{n} \omega_{k}^{h} x_{h}\) , where \(\omega_{k}\) is an n-th root of unity \((1 \leqslant k \leqslant n)\) , shows clearly how a knowledge of these n numbers entails that of the roots \(x_{k}\) , and examines in general the degree of the equation satisfied by the \(y_{k}\) ; he shows for example that if n is prime, then the \(y_{k}^{n}\) are the roots of an equation of degree n-1 whose coefficients are rational functions of a root of an equation of degree \((n-2)!\) with coefficients which may be expressed rationally by means of the coefficients of the given equation. « Here, unless I am mistaken » he concludes, « are the true principles of the solution of equations, and the analysis most appropriate to lead us there ; as we see, everything reduces to a kind of calculus of combinations, by which one finds a priori the results one should expect » ({[}11{]}, a), p.~403).

As regards the memoir of Vandermonde, though independent of that of Lagrange, it makes contact in a number of points with the latter, notably in the idea of studying rational functions of the roots taking as few distinct values as possible under permutations of the roots *, and the study of « Lagrange resolvents » which he also introduces for this purpose. His work is far from possessing the clarity and generality of that of Lagrange ; however on one point he goes decidedly further, in applying the same ideas to the equation of the division of the circle \(x^n - 1 = 0\) , for \(n\) an odd prime. While Lagrange contents himself with recalling that this equation reduces to an equation of degree \(m = (n - 1)/2\) with rational coefficients, without trying to solve it when \(n \geqslant 11\) , Vandermonde affirms that the \(m\) -th powers of the Lagrange resolvents of this latter equation are rational, by reason of the relations between the various roots of \(x^n - 1 = 0\) ; but he confines himself to verifying the wellfoundedness of this assertion for \(n = 11\) , without justifying it generally.

It is only 30 years later that the result announced by Vandermonde was completely proved by C. F. Gauss **. His decisive results on the equation \(x^n - 1 = 0\) ( \(n\) an odd prime) fit into the general programme of his memorable arithmetical researches ({[}13{]}, vol.~I, p.~413 et seq.), and they illustrate especially his mastery in the manipulation of what we now call the theory of cyclic groups. After having proved that the polynomial \(\Phi_n(x) = (x^n - 1)/(x - 1)\) is irreducible for odd prime \(n\) ***, he has the idea of writing his \(n - 1\) roots in the form \(\zeta^{g^{k}} = \zeta_{k} \quad (0 \leqslant k \leqslant n - 2)\) , where \(g\) is a primitive root of the congruence \(z^{n-1} \equiv 1 \pmod{n}\) (which, in modern terms, comes to putting in evidence the fact that the group \(\Gamma\) of the equation \(\Phi_{n}(x) = 0\) is cyclic). With every divisor \(e\) of \(n - 1\) he associates the \(f = (n - 1)/e\) « periods » \(\eta_{\nu} = \zeta_{\nu} + \zeta_{\nu + e} + \zeta_{\nu + 2e} + \cdots + \zeta_{\nu + (f - 1)e}\) ( \(1 \leqslant \nu \leqslant f\) ) and he shows in essence that the linear combinations with rational coefficients of the \(\eta_{\nu}\) form a field, generated by any one of the \(f\) periods \(\eta_{\nu}\) , and of degree \(f\) over the field of rational numbers (this field of course corresponds to the subgroup of \(\Gamma\) of order \(e\) ). We cannot here enter into the details of this analysis, and of the important arithmetic consequences which it entails; let us merely note that it gave him in particular the celebrated theorem on the possibility of constructing « by ruler and compasses » the regular polygons having a number of sides equal to a prime number of the form \(2^{2^{k}} + 1\) *. As regards the solutions by radicals of the equation \(\Phi_{n}(x) = 0\) , this follows easily from the theory of periods, applied to the \(f\) -th power of a Lagrange resolvent \(\sum_{\nu=0}^{f-1} \omega^{\nu} \eta_{\nu}\) (where \(\omega^{f} = 1\) ) **.

It is directly to Lagrange that the researches of his compatriot Ruffini are linked, contemporaneous with the Disquisitiones ; taking up the question at the point where it had been left by Lagrange, they had as their aim the proof of the impossibility of the solution « by radicals » of the « general » *** equation of the 5th degree. Ruffini's proof, prolix and obscure, remains incomplete, even though revised several times ; but it is already very close to the proof (correct in principle) which Abel obtains later ****. Its chief interest resides in the introduction of a calculus of substitutions and the first notions of the theory of groups, which

Ruffini develops to show that there exists no function of the 5 roots of the equation taking more than 2 and less than 5 values as the roots are permuted arbitrarily.

We have already said (Hist. Note to Chap. I) how this first outline of the theory of groups of permutations was developed and systematized by Cauchy some years later. But if the notions necessary for the development of Lagrange's ideas concerning substitutions are thus gradually clarified, it was still necessary to state the first principles of the theory of fields in equally clear fashion. It is this which Ruffini lacked, and which Abel and Galois were to undertake, in the last phase of the problem of the solution of algebraic equations.

Throughout his short life Abel never ceased to be preoccupied by this problem. Still almost a child, he believed to have found a formula for the solution by radicals of the general equation of the 5th degree. Later, when he recognizes his error, he does not rest until he succeeds in proving that such a formula does not exist ({[}14{]}, vol.~I, p.~166). But he does not stop there. While his competitor Jacobi develops the theory of elliptic functions as an analyst, it is the algebraic point of view which dominates Abel's work on this question, centred on the theory of equations of the division of elliptic functions ({[}14{]}, vol.~I, p.~265, 377 et passim). He thus obtains new types of equations soluble by radicals, by a method modelled on that of Gauss for the equations of the division of the circle ({[}14{]}, vol.~I, p.~310 and 358) *; a result from which he rises to the concept of « abelian » equations, of which he proves in a celebrated memoir the solubility by radicals ({[}14{]}, vol.~I, p.~478); it is on this occasion that he defines in a precise fashion the notion of an irreducible polynomial over a given field (the field generated by the coefficients of the equation he is studying) **. Finally death strikes him down in 1829, while he is attacking the general problem of the characterization of all the equations soluble by radicals, and has just communicated to Crelle and Legendre results that are already very close to those of Galois ({[}14{]}, vol.~II, p.~219-243, 269-270 and 279).

It is for the latter that it was reserved, three years later, to crown the edifice \([15]\) . Like Abel, but in even more precise fashion, he begins by defining (except for terminology) the belonging of a quantity to a field generated by given quantities, the notion of adjunction and irreducible polynomials over a given field. Given an equation \(F(x) = 0\) without multiple roots, with coefficients in a given field \(K\) , he shows successively how « it is always possible to form a function \(V\) of the roots, such that none of the values obtained by permuting the roots in all possible ways in this function is equal to another », that this function « enjoys the property that all the roots of the given equation may be expressed rationally as a function of \(V\) » and that, if \(V, V', V'', \ldots\) are the roots of the irreducible equation satisfied by \(V\) , « if \(a = f(V)\) is a root of the proposed equation, then so is \(f(V')\) » ({[}15{]}, p.~36-37); in modern terms he thus proves that \(V\) , as well as any one of its conjugates over \(K\) generates the field \(N\) of the roots of \(F\) . He then defines the group \(\Gamma\) of \(F\) as the set of permutations of the roots \(x_i\) which are obtained in substituting for \(V\) in the rational expression of each of the \(x_i\) as function of \(V\) , any one of the conjugates of \(V\) ; and he obtains at once the fundamental characterization of the elements of \(K\) by the property of being invariant under every permutation of \(\Gamma\) ({[}15{]}, p.~38-39). He then proves that if \(N\) contains the root field \(L\) of another polynomial, then the group of \(N\) over \(L\) is a normal subgroup of \(\Gamma\) (a notion which he introduces on this occasion) ({[}15{]}, p.~41 and 25-26). From this he finally deduces the criterion for the solubility of an equation by radicals, by means of reasoning which essentially runs as follows : the base field \(K\) being assumed to contain all the roots of unity, there exists by hypothesis an ascending sequence \((K_i)_{0 \leq i \leq m}\) of intermediate fields between \(K\) and \(N\) with \(K_0 = K\) , \(K_m = N\) , the field \(K_{i+1}\) being obtained by adjunction to \(K_i\) of all the roots of a binomial equation \(x^{n_i} - a_i = 0\) (with \(a_i \in K_i\) ). There exists thus in \(\Gamma\) a descending sequence \((\Gamma_i)\) of subgroups such that \(\Gamma_0 = \Gamma\) , \(\Gamma_m = \{\varepsilon\}\) (neutral element), \(\Gamma_{i+1}\) being normal in \(\Gamma_i\) and the quotient group \(\Gamma_i / \Gamma_{i+1}\) being cyclic (a case in which the group \(\Gamma\) is said to be solvable). Conversely when this is so, the use of a Lagrange resolvent shows that \(K_{i+1}\) is obtained from \(K_i\) by adjunction of all the roots of a binomial equation and therefore the equation \(F(x) = 0\) is soluble by radicals *. The impossibility of the solution by radicals of the « general » equation of degree n \textgreater{} 4 is thus a consequence of the fact that the group \(\Gamma\) of this equation, isomorphic to the symmetric group \(S_n\) (V, p.~59, Example 5), is not soluble (I, p.~142, Ex. 10 and p.~143, Ex. 16).

From the middle of the 19th century the algebraists, as we have already remarked (cf.~Hist. Note of Chap. I), enlarged considerably the domain of their investigations, until then almost entirely confined to the study of equations. In the light of Galois' discoveries it was realized that the problem of solution « by radicals » is only a particular somewhat artificial case of the general problem of classifying irrationals. It is the latter which throughout the closing years of the 19th century is attacked from various sides and numerous disparate results accumulate gradually, preparing the way for the synthesis by Steinitz.

In the first place, as far as the algebraic irrationals were concerned, a fundamental principle of classification was provided by Galois theory, which reduced the study of an algebraic equation to that of its group. For this reason it is above all the theory of permutation groups, of which it is not our task to speak here (cf.~Hist. Note to Chap. I), which in pure Algebra formed the principal object of research of that period. The remaining progress in the theory of algebraic fields was due to the development, at the same time, of the Theory of Numbers and Algebraic Geometry. Actually this progress concerned mainly the exposition of the theory, and was mostly due to Dedekind {[}18{]}, who introduced the notions of field and ring *, and (in connexion with his researches on hypercomplex systems) develops systematically the linear aspect of the theory of extensions ({[}18{]}, vol.~3, p.~33 et seq.). It is he also who considers the Galois group as consisting of automorphisms of the extension being considered, and not merely as a group of permutations of the roots of an equation, and he proves (for number fields) his fundamental theorem on the linear independence of automorphisms ({[}18{]}, vol.~3, p.~29) as well as the existence of normal bases of a Galois extension ({[}18{]}, vol.~2, p.~433). Finally he tackles the problem of algebraic extensions of infinite degree and notes that Galois theory cannot be applied as it stands (an arbitrary subgroup of a Galois group need not always be identical with the group of the extension with respect to a subextension), and by a daring intuition he already tries to contemplate the Galois group as a topological group ** --- an idea which reaches maturity only with the theory of Galois extensions of infinite degree, developed by Krull in 1928 {[}24{]}.

Parallel to this evolution the notion of a transcendental element over a field is made precise. The existence of transcendental numbers is proved for the first time by Liouville in 1844, by an explicit process of construction, based on the theory of diophantine approximations {[}17{]}; Cantor in 1874 gave another « non-constructive » proof, using simple reflexions on the power of sets (cf.~V, p.~147), Ex. 3); finally Hermite proved in 1873 the transcendence of e and Lindemann in 1882 that of \(\pi\) by a method analogous to Hermite's, thus finally closing the discussion on the antique problem of the quadrature of the circle*.

Regarding the role of transcendental numbers in algebraic calculations, Kronecker observes in 1882 that if x is transcendental over a field K, then the field \(\mathrm{K}(x)\) is isomorphic to the field of rational fractions \(\mathrm{K}(\mathrm{X})\) ({[}19{]}, a), p.~7). He actually makes the adjunction of indeterminates to a field the corner stone of his exposition of the theory of algebraic numbers ({[}19{]}, a)). On the other hand, Dedekind and Weber show in the same year ({[}18{]}, vol.~I, p.~238) how arithmetic methods may be used to set up a theory of algebraic curves. The analogies between Arithmetic and Algebraic Geometry, which will turn out to be so fruitful for the one and the other are thus seen to appear from several directions.

In all these researches the fields which occur consist of « concrete » elements in the sense of classical mathematics --- (complex) numbers or functions of complex variables **. But already Kronecker in 1882 was perfectly aware of the fact (only obscurely glimpsed by Gauss and Galois) that the « indeterminates » play in his theory only the role of basis elements of an algebra and not that of variables in the sense of Analysis ({[}19{]}, a), p.~93-95); and in 1887 he develops this idea, combined with a vast programme which aims at nothing less than to re-lay the foundations for all of mathematics by rejecting everything that cannot be reduced to algebraic operations on integers (cf.~Hist. Note to Book I, Chap. IV). It is on this occasion that, taking up an idea of Cauchy {[}16{]} who had defined the field C of complex numbers as the residue class field \(\mathbf{R}[\mathbf{X}] / (\mathbf{X}^2 + 1)\) , Kronecker shows how the theory of algebraic numbers is entirely independent of the « fundamental theorem of algebra » and even of the theory of real numbers, since every algebraic number field (of finite degree) is isomorphic to a residue class field \(\mathbf{Q}[\mathbf{X}] / (f)\) (where \(f\) is an irreducible polynomial over \(\mathbf{Q}\) ) {[}19{]}, b). As H. Weber {[}20{]} remarks some years later in developing a first sketch of an axiomatic theory of fields, this method of Kronecker's applies in fact to any base field K. Weber indicates in particular that a field \(\mathbf{Z}/(p)\) ( \(p\) prime) may be taken for K, allowing thus the calculus of congruences « modulo \(p\) » to be encompassed by the theory of fields; the former had taken shape in the second half of the 18th century at the hands of Euler, Lagrange, Legendre and Gauss, and the analogy it presented with the theory of algebraic equations had been readily observed; in developing that analogy, Galois (in the light of his researches in group theory) had not hesitated to introduce « ideal roots » of an irreducible congruence modulo \(p\) * , and had indicated its chief properties ({[}15{]}, p.~15-23) **. When Kronecker's method is applied to \(\mathbf{Z}/(p)\) , the result is (except for terminology) the presentation which Serret and Dedekind ({[}18{]}, vol.~I, p.~40) had already given of these « Galois imaginaries ».

To all these examples of « abstract fields » there is added, at the turn of the century, the field of formal power series introduced by Veronese {[}21{]} and above all the \(p\) -adic fields of Hensel {[}22{]}. It is the discovery of the latter which leads Steinitz (as he says explicitly) to isolate the abstract notions common to all these theories, in a basic paper {[}23{]} which may be considered as having given rise to the current conception of Algebra. Developing systematically the consequences of the axioms of commutative fields, he thus introduces the notions of prime field, of separable (algebraic) elements, of perfect field, he defines the transcendence degree of an extension and finally proves the existence of algebraically closed extensions of an arbitrary field.

More recently the theory of Steinitz has been completed in several important respects. On the one hand the work of Artin has put in evidence the linear character of Galois theory \([25]\) . On the other hand, the general notion of derivation (modelled on the formal properties of classical Differential Calculus) sensed by Dedekind ({[}18{]}, vol.~2, p.~412), introduced by Steinitz in the particular case of a field of rational fractions ({[}23{]}, p.~209-212), has been used with success in the study (essential for modern Algebraic Geometry) of transcendental extensions, and notably in the generalization of the notion of separability to the latter \([26]\) .

